\subsubsection{4.1 研究方案}

xxx

% 方案一

% 方案二

% A
\textbf{A. 面向端地协同的频谱数据存传耦合价值评估方法}

本项目拟研究面向端地协同的频谱数据存传耦合价值评估方法，整体框架如图\ref{fig:3A}所示。

针对无人机端侧算力、存储和通信带宽受限，宽频段频谱数据中背景冗余高、关键事件证据易丢失的问题，本研究提出\textbf{面向事件证据窗口的频谱数据多维价值表征方法}，将连续频谱流组织为以频谱事件为核心的证据单元，由事件主体、触发前背景、触发后恢复过程、频域邻域和任务语境共同刻画信号的出现、持续、变化及其任务关联关系。在此基础上，构建端侧频谱知识库，维护任务频段、历史背景基线、常见干扰模板、已知信号模式、已复核事件原型和地面端反馈规则等轻量信息，为事件价值判断提供先验约束和历史参照。由此，从电磁异常性、任务相关性、历史稀缺性、证据完整性和时效敏感性等维度刻画频谱数据价值，使频谱片段的重要性不仅取决于当前信号强度，还取决于其相对于背景是否异常、相对于历史是否少见、相对于任务是否关键以及是否具备后端复核所需的完整证据。

针对端侧无法进行全量深度分析、存储保留与链路回传需求相互影响的问题，本研究进一步提出\textbf{存传耦合驱动的轻量化频谱价值计算方法}。端侧在滑动频谱流上维护背景基线、频段占用、事件频次、历史稀缺度和任务频段命中情况等增量状态，将价值计算由全量分析转化为局部更新；地面端则承担复杂特征分析、相似事件检索和知识库更新，并将结果以轻量参数或模板形式回写端侧，从而形成端侧增量初算与地面复杂辅算相结合的计算模式。在此基础上，分别计算频谱事件的存储价值和传输价值：存储价值衡量事件是否需要在端侧保留更完整的原始证据或时频窗口，传输价值衡量事件是否需要立即占用链路回传地面端。二者通过事件生命周期状态、链路可用性、本地存储压力和地面补证需求形成耦合关系，使端侧能够根据事件是否已回传、是否被确认、是否需要补传以及当前链路和存储条件，动态调整本地留存和回传优先级。为提高端侧计算效率，本研究将连续复杂决策转化为有限动作集合上的边际收益比较，在摘要保存、ROI 保存、I/Q 窗口暂存、即时回传、延迟补传和本地淘汰等动作之间进行快速选择：

\begin{equation}
	u_e^{*} = \arg\max_{u\in\mathcal{A}_{lite}} \Delta Q(e,u),
\end{equation}
\begin{equation}\notag
	\Delta Q(e,u)
	= \widehat{R}_{\theta}(e,u)
	- \lambda_s \Delta C_s(u)
	- \lambda_b \Delta C_b(u)
	- \lambda_c \Delta C_c(u),
\end{equation}


其中，$\mathcal{A}_{lite}$ 表示端侧可执行的轻量动作集合，$\widehat{R}_{\theta}(e,u)$ 表示动作 $u$ 对后端复核、实时告警和任务执行的预期收益，$\Delta C_s(u)$、$\Delta C_b(u)$ 和 $\Delta C_c(u)$ 分别表示新增存储、带宽和计算开销。由于候选动作集合规模有限，端侧仅需进行低维边际收益计算即可完成存传协同决策，从而在无独立 AI 算力条件下实现高效价值评估。

针对动态电磁环境中背景噪声、干扰模式、任务目标和链路条件持续变化导致固定价值规则易失效的问题，本研究进一步提出\textbf{端地协同反馈驱动的频谱价值动态校准机制}。地面端在接收端侧回传的摘要、ROI 或 I/Q 证据后，利用更强算力进行复杂特征提取、跨任务相似事件检索、干扰模板匹配和人工复核，形成面向端侧更新的结构化反馈。基于地面端反馈，将端侧价值模型更新表示为知识库约束下的在线校准问题：
\begin{equation}
	\begin{aligned}
		(\theta_{t+1},\mathcal{K}_{t+1})
		= \arg\min_{\theta,\mathcal{K}}
		\sum_{e\in\mathcal{B}_t}
		\ell\!\left(
		\mathcal{F}_{\theta}
		(\mathbf{z}_e,\mathbf{h}_t,\mathcal{K},\mathbf{s}_t,a_e),
		y_e
		\right) 
		+ \Omega(\theta-\theta_t)
		+ \gamma D(\mathcal{K},\mathcal{K}_t),
	\end{aligned}
\end{equation}
其中，$\mathcal{B}_t$ 表示近期被地面端复核的事件集合，$y_e$ 表示结构化反馈，$\Omega(\cdot)$ 用于约束端侧轻量模型更新幅度，$D(\cdot)$ 用于约束知识库更新代价。同时，引入场景门控机制，使模型参数随链路质量、剩余存储、电磁活跃程度和任务阶段自适应调整：链路受限时提高本地证据保留倾向，存储紧张时提高压缩与淘汰强度，重点侦察阶段提高任务频段和未知事件的处理优先级。通过该机制，形成“端侧增量计算--地面辅助分析--结构化反馈校准--知识库更新”的闭环，使频谱数据价值模型能够随任务、链路和电磁环境变化持续演化。

综上，本方案构建了端地协同的频谱数据存传耦合价值评估框架，通过事件证据窗口和端侧知识库实现多维价值表征，通过有限动作边际收益比较实现轻量化存传价值计算，并通过地面端反馈与知识库约束更新实现动态校准，为后续频谱数据语义压缩、分级存储、低带宽回传和可信管理提供基础支撑。
% B

% C
\textbf{C. 链上链下协同的频谱数据可信管理方法}

本项目拟研究\textbf{链上链下协同的频谱数据可信管理方法}，重点突破基于可信硬件的数据验证、基于密码累加器的动态摘要维护以及基于可验证索引的可信查询三个关键环节。

端侧频谱数据库中频谱数据记录规模大、校验频次高，若仅依赖数据库本地日志或中心化校验结果，难以保证验证过程本身可信。为解决这一问题，本课题研究\textbf{基于可信硬件的数据验证方法}，将链下频谱数据的校验逻辑嵌入可信执行环境。首先在可信硬件内部对待验证数据 $d_i$ 提取由数据版本、状态标识和数据库记录状态构成的状态描述 $s_i$，并生成对应的验证摘要：
\[
h_i = H(d_i \,\|\, s_i),
\]
其中，$H(\cdot)$ 表示安全哈希函数，$h_i$ 表示链下频谱数据状态的摘要结果。随后，由可信执行环境对摘要结果和验证过程进行签名或证明，形成可上链记录的可信验证信息：
\[
r_i = \mathrm{Sign}_{\mathrm{TEE}}(h_i, t_i),
\]
其中，$t_i$ 表示验证时间，$r_i$ 表示由可信硬件生成的验证记录。链上仅记录 $h_i$ 与 $r_i$ 等轻量化可信信息，不直接保存大规模频谱数据。后续在数据状态更新或查询校验过程中，可重新计算链下数据摘要并与链上记录进行比对：
\[
V_{\mathrm{TEE}}(d_i, s_i, h_i, r_i) = 1,
\]
当且仅当 $H(d_i \,\|\, s_i) = h_i$ 且 $\mathrm{Verify}_{\mathrm{TEE}}(r_i) = 1$ 时成立，其中 $V_{\mathrm{TEE}}(\cdot)$ 表示基于可信硬件的数据验证函数，$\mathrm{Verify}_{\mathrm{TEE}}(\cdot)$ 表示对可信硬件验证记录的校验过程。综上，本课题通过可信执行环境对链下频谱数据校验逻辑进行隔离执行，并将验证摘要、验证时间等信息绑定生成可复核的验证记录，再将轻量化验证信息写入链上。后续数据状态维护和查询校验过程中，可通过可信硬件验证记录与链上摘要信息对链下数据状态进行一致性比对，从而为链上链下协同管理提供可靠验证依据。

频谱数据在端侧持续产生新增、修改和删除等操作，若每次数据状态变化后都对链下数据库进行全量校验，将难以满足低开销和高时效的管理要求。为解决这一问题，本课题研究\textbf{基于密码累加器的动态摘要机制}，将频谱数据变更抽象为可增量处理的更新事件，并利用密码累加器维护链下数据状态对应的动态摘要。首先，定义频谱数据在时刻 $t$ 的变更事件为 $e_t = (\mathrm{op}_t, \mathrm{id}_t, \Delta d_t)$，其中 $\mathrm{op}_t$ 表示变更类型，$\mathrm{id}_t$ 表示频谱数据记录标识，$\Delta d_t$ 表示变更内容。针对新增操作，将新增记录的摘要元素加入累加器；针对删除操作，将对应记录的摘要元素从累加器状态中移除；针对修改操作，将其转化为旧摘要元素删除与新摘要元素加入的组合过程，从而将不同类型的数据变更统一纳入动态摘要维护框架。给定当前摘要状态 $A_t$，通过累加器更新函数得到变更后的摘要状态：
\[
A_{t+1} = U_{\mathrm{Acc}}(A_t, e_t),
\]
其中，$U_{\mathrm{Acc}}(\cdot)$ 表示动态摘要更新函数，$A_{t+1}$ 表示频谱数据变更后的摘要状态。其次，针对变更后的数据状态生成增量验证凭据 $\pi_t$，并利用累加器验证函数判断变更结果是否已被正确纳入当前摘要状态：
\[
V_{\mathrm{Acc}}(A_{t+1}, e_t, \pi_t) = 1,
\]
其中，$V_{\mathrm{Acc}}(\cdot)$ 表示累加器验证函数。通过该机制，系统只需校验受变更影响的频谱数据记录及其摘要关系，而不需要对链下数据库中的全部频谱数据进行重复校验。为支撑链上链下状态同步维护，将更新后的摘要状态 $A_{t+1}$ 及其对应的变更标识写入链上，使链上摘要信息能够反映链下频谱数据状态的最新变化。后续在数据查询或状态核验过程中，可根据链上摘要状态与链下数据记录重新生成的摘要关系进行比对，判断链下数据是否与当前可信状态一致。综上，本课题利用密码累加器的增量验证特性，将频谱数据更新后的状态维护由全量校验转化为局部增量校验，实现链下数据变更状态与链上摘要信息的同步维护，降低频谱数据高频更新条件下的状态维护开销。

海杂波场景下，高价值频谱数据往往分散在不同观测节点、不同频段和不同时段的频谱记录中，仅依赖传统数据库索引进行查询，难以同时满足复杂条件下的快速定位与结果校验需求。为解决这一问题，本课题研究\textbf{基于可验证索引的频谱数据可信查询机制}，通过分析频谱数据的关联特征构建可验证索引模型，将查询定位过程与结果验证过程统一到同一索引框架中。首先，构建面向复杂频谱查询的索引建模方法。对链下频谱数据记录 $d_i$ 提取用于查询定位的关联特征 $\phi_i$，并将其与前述可信验证环节生成的数据摘要 $h_i$ 进行绑定，形成可验证索引项：
\[
z_i = (\phi_i, h_i, \mathrm{id}_i),
\]
其中，$\phi_i$ 表示频谱数据记录的关联特征，$h_i$ 表示该记录对应的可信摘要，$\mathrm{id}_i$ 表示频谱数据记录标识。在此基础上，构建索引模型 $I(\cdot)$，使查询过程不仅能够依据关联特征定位候选数据，还能够保留与数据摘要相关的校验依据。给定查询请求 $q$，通过索引匹配函数获得候选结果集合：
\[
R_q = I(q),
\]
其中，$R_q$ 表示满足查询条件的候选频谱数据结果。其次，设计面向复杂频谱查询的索引编码与验证信息生成方法。针对查询过程中可能涉及多条件组合、范围筛选和关联检索等操作，根据查询请求 $q$、候选结果集合 $R_q$ 以及对应的索引访问路径生成查询验证信息：
\[
P_q = \mathrm{Prove}_I(q, R_q),
\]
其中，$P_q$ 表示与查询路径和返回结果相关的验证信息。该验证信息用于描述查询结果如何由索引模型得到，以及返回结果对应的摘要关系是否能够与当前链上摘要状态相匹配。通过该设计，查询过程不再仅返回数据结果，还同步生成可用于复核的验证依据。最后，构建查询结果校验方法，将查询请求、返回结果、验证信息与当前链上摘要状态进行一致性校验：
\[
V_Q(q, R_q, P_q, A_t) = 1,
\]
当且仅当 $\mathrm{Check}_{\mathrm{path}}(q, R_q, P_q) = 1$ 且 $\mathrm{Check}_{\mathrm{state}}(R_q, P_q, A_t) = 1$ 时成立，其中 $V_Q(\cdot)$ 表示查询结果校验函数，$A_t$ 表示当前链上摘要状态，$\mathrm{Check}_{\mathrm{path}}(\cdot)$ 用于校验查询路径与返回结果之间的一致性，$\mathrm{Check}_{\mathrm{state}}(\cdot)$ 用于校验返回结果与链上摘要状态之间的一致性。当校验通过时，说明查询结果由当前索引路径正确返回，且与链上可信状态一致；当校验失败时，说明查询路径、返回结果或摘要状态之间存在不匹配。综上，本课题通过基于可验证索引的频谱数据可信查询机制，将复杂频谱查询、索引路径复核和结果状态校验结合起来，提升高价值频谱数据的查询效率与结果可验证性，为海杂波场景电磁频谱数据的高可信管理提供技术支撑。

% 方案三